\section*{Расчёт цен покупателя и продавца европейских опционов в однопериодической триномиальной модели}

\subsection*{1. Исходные параметры}

Пусть:
\[
S_0 = 190, \quad K = 190, \quad T = 0.25
\]
\[
r = 5\% \Rightarrow R = 1 + rT = 1.0125
\]

Параметры триномиальной модели:
\[
u = 1.1, \quad m = 1.0, \quad d = 0.9
\]

Тогда возможные значения цены актива:
\[
S_u = 209, \quad S_m = 190, \quad S_d = 171
\]

\subsection*{2. Риск-нейтральные вероятности}

Вероятности удовлетворяют условиям:
\[
p_u + p_m + p_d = 1
\]
\[
1.1p_u + 1.0p_m + 0.9p_d = 1.0125
\]

Выберем:
\[
p_u = 0.30, \quad p_m = 0.50, \quad p_d = 0.20
\]

\subsection*{3. Цена опциона Call}

Выплаты:
\[
C_u = \max(209 - 190, 0) = 19
\]
\[
C_m = 0, \quad C_d = 0
\]

Математическое ожидание:
\[
E[C_T] = 0.30 \cdot 19 = 5.7
\]

Цена опциона:
\[
C = \frac{E[C_T]}{R} = \frac{5.7}{1.0125} \approx 5.63
\]

\subsection*{4. Цена опциона Put}

Выплаты:
\[
P_u = 0, \quad P_m = 0
\]
\[
P_d = \max(190 - 171, 0) = 19
\]

Математическое ожидание:
\[
E[P_T] = 0.20 \cdot 19 = 3.8
\]

Цена опциона:
\[
P = \frac{E[P_T]}{R} = \frac{3.8}{1.0125} \approx 3.75
\]

\subsection*{5. Цены покупателя и продавца}

В идеальной модели:
\[
C_{bid} = C_{ask} = 5.63, \quad
P_{bid} = P_{ask} = 3.75
\]

С учётом спрэда (примерно 5\%):
\[
C_{bid} \approx 5.50, \quad C_{ask} \approx 5.75
\]
\[
P_{bid} \approx 3.65, \quad P_{ask} \approx 3.85
\]

\subsection*{6. Вывод}

Полученные цены ниже рыночных, так как модель:
\begin{itemize}
\item однопериодическая,
\item использует заниженную волатильность,
\item не учитывает реальные рыночные факторы.
\end{itemize}
